\documentclass[a4paper,twoside]{article}

\usepackage{morewrites}

\usepackage[svgnames,dvipsnames,x11names]{xcolor}
\usepackage{graphicx}
\usepackage[english]{babel}
\usepackage[colorlinks=true, linkcolor=NavyBlue, urlcolor=NavyBlue, citecolor=NavyBlue]{hyperref}
\usepackage[a4paper]{geometry}
\usepackage{ts-typesetting}
\usepackage{ts-fonts}
\usepackage{float}
\usepackage{xcolor}
\usepackage{epigraph}
\usepackage{marginnote}
\usepackage{multicol}
\usepackage{progressbar}
\usepackage{graphicx}
\usepackage{tcolorbox}
\usepackage{fvextra}
\usepackage{amsmath}
\usepackage{seqsplit}
\renewcommand*{\marginfont}{
\footnotesize
\sloppy
\emergencystretch=1em
\hyphenpenalty=50
\exhyphenpenalty=50
}
\newlength{\tsmarginparavail}
\newlength{\tsmarginparalt}
\newlength{\tsmarginparbuf}
\setlength{\tsmarginparbuf}{6mm}
\makeatletter
\AtBeginDocument{
\setlength{\tsmarginparavail}{\dimexpr\paperwidth-\textwidth-1in-\oddsidemargin-\marginparsep-\tsmarginparbuf\relax}
\ifdim\tsmarginparavail<\z@\setlength{\tsmarginparavail}{\z@}\fi
\if@twoside
\setlength{\tsmarginparalt}{\dimexpr1in+\evensidemargin-\marginparsep-\tsmarginparbuf\relax}
\ifdim\tsmarginparalt<\z@\setlength{\tsmarginparalt}{\z@}\fi
\ifdim\tsmarginparalt<\tsmarginparavail
\setlength{\tsmarginparavail}{\tsmarginparalt}
\fi
\fi
\ifdim\tsmarginparavail<\marginparwidth
\setlength{\marginparwidth}{\tsmarginparavail}
\fi
}
\makeatother
\usepackage{ts-callouts}
\usepackage{ts-code}

\usepackage{lastpage}
\usepackage{refcount}
\makeatletter
\def\ps@plain{
\let\@oddhead\@empty
\let\@evenhead\@empty
\def\@oddfoot{
\hfil
\ifnum\getpagerefnumber{LastPage}>1\relax
\thepage
\fi
\hfil}
\let\@evenfoot\@oddfoot
}
\makeatother

\title{Admonitions}
\author{}\date{}
\begin{document}
\maketitle
Admonitions --- callouts --- are little blocks for surfacing notes, warnings, tips,
and whatever else you need to highlight. In TMark they are \textbf{containers}, and
the canonical spelling is the \tscodeinline{:::} fence:
\begin{tscode}[lang=md, engine=pygments]
\begin{Verbatim}[commandchars=\\\{\},breaklines, breakanywhere, commandchars=\\\{\}]
::: note \PYZob{}title=\PYZdq{}This is a Note\PYZdq{}\PYZcb{}
Any number of other Markdown elements.

This is the second paragraph.
:::
\end{Verbatim}
\end{tscode}
\begin{tscallout}[kind=note, title={This is a Note}]
Any number of other indented Markdown elements.

This is the second paragraph.
\end{tscallout}
Folding is an attribute, not a second fence family:
\begin{tscode}[lang=md, engine=pygments]
\begin{Verbatim}[commandchars=\\\{\},breaklines, breakanywhere, commandchars=\\\{\}]
::: note \PYZob{}title=\PYZdq{}This is a Note\PYZdq{} collapsed=true\PYZcb{}
Any number of other Markdown elements.

This is the second paragraph.
:::
\end{Verbatim}
\end{tscode}
\begin{tscallout}[kind=note, title={This is a Note}, collapsed]
Any number of other indented markdown elements.

This is the second paragraph.
\end{tscallout}
\begin{tscallout}[kind=info, title={The PyMdownX spellings are kept indefinitely}]
\tscodeinline{!!! note "Title"} and \tscodeinline{??? note "Title"} (collapsed) / \tscodeinline{???+} (expanded)
remain accepted sugar for callouts, because MkDocs Material renders them
natively. They are one node whatever the spelling. Use whichever suits the
document; \tscodeinline{:::} is what the printer emits.
\end{tscallout}
Print has no folding, so the strategy is declared once:
\begin{tscode}[lang=yaml, engine=pygments]
\begin{Verbatim}[commandchars=\\\{\},breaklines, breakanywhere, commandchars=\\\{\}]
\PY{n+nt}{press}\PY{p}{:}
\PY{+w}{  }\PY{n+nt}{details}\PY{p}{:}\PY{+w}{ }\PY{l+lScalar+lScalarPlain}{expand}\PY{+w}{      }\PY{c+c1}{\PYZsh{} expand (default) | reference}
\end{Verbatim}
\end{tscode}
\tscodeinline{reference} moves the body of a collapsed callout to a grouped end-section
(``Solutions'', ``Warnings'', ...) and leaves a ``See page N'' link in its place.
\section{\LaTeX{} Rendering}
TeXSmith maps callouts onto the \tscodeinline{tcolorbox} package through the \tscodeinline{ts-\allowbreak{}callouts}
fragment, which provides the \tscodeinline{tscallout} environment:
\begin{tscode}[lang=latex, engine=pygments]
\begin{Verbatim}[commandchars=\\\{\},breaklines, breakanywhere, commandchars=\\\{\}]
\PY{k}{\PYZbs{}begin}\PY{n+nb}{\PYZob{}}tscallout\PY{n+nb}{\PYZcb{}}[kind=note, title=\PY{n+nb}{\PYZob{}}This is a Note\PY{n+nb}{\PYZcb{}}]
…
\PY{k}{\PYZbs{}end}\PY{n+nb}{\PYZob{}}tscallout\PY{n+nb}{\PYZcb{}}
\end{Verbatim}
\end{tscode}
Template authors restyle them by redefining that environment or by appending to
its \tscodeinline{pgfkeys} family. Set the global look with \tscodeinline{press.callouts.style} (or via
\tscodeinline{-\allowbreak{}-\allowbreak{}attribute press.callouts.style=<style>}):
\begin{tscode}[lang=yaml, engine=pygments]
\begin{Verbatim}[commandchars=\\\{\},breaklines, breakanywhere, commandchars=\\\{\}]
\PY{n+nn}{\PYZhy{}\PYZhy{}\PYZhy{}}
\PY{n+nt}{press}\PY{p}{:}
\PY{+w}{  }\PY{n+nt}{callouts}\PY{p}{:}
\PY{+w}{    }\PY{n+nt}{style}\PY{p}{:}\PY{+w}{ }\PY{l+lScalar+lScalarPlain}{classic}\PY{+w}{  }\PY{c+c1}{\PYZsh{} fancy | classic | minimal}
\PY{n+nn}{\PYZhy{}\PYZhy{}\PYZhy{}}
\end{Verbatim}
\end{tscode}
\begin{itemize}
\item \tscodeinline{fancy} (default): colored headings with icons.
\item \tscodeinline{classic}: black-and-white layout with a bold left rule.
\item \tscodeinline{minimal}: subtle border, rounded corners, and no icons.
\end{itemize}
\tscodeinline{press.callout\_style} and \tscodeinline{press.admonition\_style} are the deprecated names of
that key.
\begin{tsdiv}{tab}[title={Fancy Admonitions}]
\begin{figure}[H]
\centering
\includegraphics[width=0.7\linewidth]{snippet-<HASH>.pdf}
\caption[Demo]{Demo}
\end{figure}
\end{tsdiv}
\begin{tsdiv}{tab}[title={Classic Admonitions}]
\begin{figure}[H]
\centering
\includegraphics[width=0.7\linewidth]{snippet-<HASH>.pdf}
\caption[Demo]{Demo}
\end{figure}
\end{tsdiv}
\begin{tsdiv}{tab}[title={Minimal Admonitions}]
\begin{figure}[H]
\centering
\includegraphics[width=0.7\linewidth]{snippet-<HASH>.pdf}
\caption[Demo]{Demo}
\end{figure}
\end{tsdiv}
\section{Custom types}
A new \emph{kind} of callout is a declaration, not a new container name. Declare what
it \textbf{is} under \tscodeinline{press.declare}, and how it \textbf{looks} under \tscodeinline{press.callouts}:
\begin{tscode}[lang=yaml, engine=pygments]
\begin{Verbatim}[commandchars=\\\{\},breaklines, breakanywhere, commandchars=\\\{\}]
\PY{n+nt}{press}\PY{p}{:}
\PY{+w}{  }\PY{n+nt}{declare}\PY{p}{:}
\PY{+w}{    }\PY{n+nt}{admonitions}\PY{p}{:}
\PY{+w}{      }\PY{n+nt}{solution}\PY{p}{:}
\PY{+w}{        }\PY{n+nt}{name}\PY{p}{:}\PY{+w}{ }\PY{l+lScalar+lScalarPlain}{Solution}
\PY{+w}{        }\PY{n+nt}{group}\PY{p}{:}\PY{+w}{ }\PY{l+lScalar+lScalarPlain}{Solutions}\PY{+w}{                          }\PY{c+c1}{\PYZsh{} section title under `reference`}
\PY{+w}{        }\PY{n+nt}{reference}\PY{p}{:}\PY{+w}{ }\PY{l+s}{\PYZdq{}}\PY{l+s}{See}\PY{n+nv}{ }\PY{l+s}{page}\PY{n+nv}{ }\PY{l+s}{\PYZob{}page\PYZcb{}}\PY{n+nv}{ }\PY{l+s}{for}\PY{n+nv}{ }\PY{l+s}{the}\PY{n+nv}{ }\PY{l+s}{solution}\PY{l+s}{\PYZdq{}}
\PY{+w}{  }\PY{n+nt}{callouts}\PY{p}{:}
\PY{+w}{    }\PY{n+nt}{solution}\PY{p}{:}\PY{+w}{ }\PY{p+pIndicator}{\PYZob{}}\PY{n+nt}{icon}\PY{p}{:}\PY{+w}{ }\PY{l+s}{\PYZdq{}}\PY{l+s}{🎓}\PY{l+s}{\PYZdq{}}\PY{p+pIndicator}{,}\PY{n+nt}{ color}\PY{p}{:}\PY{+w}{ }\PY{l+s}{\PYZdq{}}\PY{l+s}{\PYZsh{}123456}\PY{l+s}{\PYZdq{}}\PY{p+pIndicator}{\PYZcb{}}\PY{+w}{      }\PY{c+c1}{\PYZsh{} quote it: a bare \PYZsh{} starts a YAML comment}
\end{Verbatim}
\end{tscode}
\begin{tscode}[lang=md, engine=pygments]
\begin{Verbatim}[commandchars=\\\{\},breaklines, breakanywhere, commandchars=\\\{\}]
\PYZhy{}\PYZhy{}\PYZhy{}
press:
  declare:
    admonitions:
\PY{g+gu}{      solution: \PYZob{}name: Solution, group: Solutions\PYZcb{}}
\PY{g+gu}{\PYZhy{}\PYZhy{}\PYZhy{}}

::: solution \PYZob{}title=\PYZdq{}Exercise 3\PYZdq{}\PYZcb{}
Apply the chain rule twice.
:::
\end{Verbatim}
\end{tscode}
Theorem environments are exactly this, with a counter attached --- see
\href{notes.md\#theorems}{Theorems}.
\section{Built-in Admonition Types}
The following types are built in:
\begin{tscallout}[kind=note]
A Note
\end{tscallout}
\begin{tscallout}[kind=tip]
A Tip
\end{tscallout}
\begin{tscallout}[kind=warning]
A Warning
\end{tscallout}
\begin{tscallout}[kind=important]
An important notice
\end{tscallout}
\begin{tscallout}[kind=danger]
A danger notice
\end{tscallout}
\begin{tscallout}[kind=info]
An info notice
\end{tscallout}
\begin{tscallout}[kind=hint]
A Hint
\end{tscallout}
\begin{tscallout}[kind=seealso]
A see-also notice
\end{tscallout}
\begin{tscallout}[kind=question]
A question notice
\end{tscallout}
\begin{tscallout}[kind=abstract]
An abstract
\end{tscallout}
\section{Content tabs}
A \tscodeinline{tabs} container holds \tscodeinline{tab} containers, each with a \tscodeinline{title=}:
\begin{tscode}[lang=md, engine=pygments]
\begin{Verbatim}[commandchars=\\\{\},breaklines, breakanywhere, commandchars=\\\{\}]
:::: tabs
::: tab \PYZob{}title=Windows\PYZcb{}
Windows is a Microsoft operating system.
:::
::: tab \PYZob{}title=Linux\PYZcb{}
Linux is an open\PYZhy{}source operating system.
:::
::::
\end{Verbatim}
\end{tscode}
On the web the reader sees one at a time; print has no interaction, so the
paged writers render the tabs in sequence, each as a titled block. The
PyMdownX spelling \tscodeinline{=== "Windows"} followed by its four-space-indented body is
accepted indefinitely --- MkDocs Material renders it natively, and this very page
uses it above --- but \tscodeinline{:::} is the canonical form and what the printer emits.
\section{Other containers}
The container names TMark knows form a closed registry. Besides the callout
types, \tscodeinline{aside}, \tscodeinline{figure}, \tscodeinline{tabs} and \tscodeinline{tab}, there are two layout containers:
\begin{tscode}[lang=md, engine=pygments]
\begin{Verbatim}[commandchars=\\\{\},breaklines, breakanywhere, commandchars=\\\{\}]
::: multicolumn \PYZob{}cols=2\PYZcb{}
The content flows in two columns.
:::

::: div \PYZob{}.grid .cards\PYZcb{}
A container that means nothing: a hook for classes and an id.
:::
\end{Verbatim}
\end{tscode}
A \tscodeinline{::: name} whose name is unknown raises \tscodeinline{container-\allowbreak{}unknown} and renders its
content transparently rather than silently becoming a \tscodeinline{<div>}.
\end{document}
